\section{Optimal profiles and smooth source realization}
\label{sec:construction}

\subsection{Explicit trial profiles}

The exact optimizer in \cref{thm:main} is the top eigenfunction of
$\mathcal K_{y/2}$.  Two elementary profiles give the explicit lower bounds in
\cref{eq:global-explicit-bracket} and make both asymptotic regimes concrete.
For the vacuum-controlled regime, set $L=yR$, choose $q=0$, and use
\begin{equation}
 p_L(x)=\sqrt{\frac{3}{L^3}}\,x\,\mathbf1_{(0,L)}(x).
 \label{eq:linear-profile}
\end{equation}
The vacuum part of the thermal form is
\begin{align}
 \ip{p_L}{h_0^{-1}p_L}
 &=\frac{3}{\pi L^3}
 \int_0^L\!\int_0^L xy
 \log\frac{x+y}{\abs{x-y}}\dd x\dd y \notag\\
 &=\frac{3L}{2\pi}.
 \label{eq:linear-profile-covariance}
\end{align}
For the second equality, scale $x=Lu$, $y=Lv$, split the unit square across
the diagonal, and use
\begin{equation}
 \int_0^1 t\log\frac{1+t}{1-t}\dd t=1.
\end{equation}
Since $\coth(\beta h_0/2)\geq1$, the full thermal dephasing exponent is at
least the value in \cref{eq:linear-profile-covariance}.  On the other hand,
$\EK=1/2$ for the normalized datum.  Thus
$\Gamma/(\EK R)\geq3L/(\pi R)=3y/\pi$.

For the thermal regime, use instead the normalized first eigenfunction of the
Dirichlet--Neumann Green operator,
\begin{equation}
 p_L^{\rm th}(x)=\sqrt{\frac2L}\sin\frac{\pi x}{2L}
 \mathbf1_{(0,L)}(x).
 \label{eq:thermal-profile}
\end{equation}
Since $\coth z\geq z^{-1}$ and
$\ip{p_L^{\rm th}}{h_0^{-2}p_L^{\rm th}}=4L^2/\pi^2$, this datum gives
\begin{equation}
 \frac{\Gamma}{\EK R}\geq\frac{8y^2}{\pi^3}.
\end{equation}
The first profile and the vacuum Schur bound also show
\begin{equation}
 \frac{3}{\pi}
 \leq\liminf_{y\downarrow0}\frac{C_{\rm opt}(y)}y
 \leq\limsup_{y\downarrow0}\frac{C_{\rm opt}(y)}y
 \leq\frac{4\operatorname{arsinh}(1)}\pi.
 \label{eq:small-support-gap}
\end{equation}
The ratio of the analytic upper and lower linear coefficients is
$4\operatorname{arsinh}(1)/3\simeq1.17516$.

For a target exponent $\Gamma_\star$, scaling
\cref{eq:linear-profile} by an amplitude $A$ yields the sufficient energy
bound
\begin{equation}
 \EK\leq\frac{\pi\Gamma_\star}{3L}.
 \label{eq:constructive-energy}
\end{equation}
This is an upper bound on the energy needed by this simple family.  The sharp
half-line momentum infimum at inverse temperature $\beta$ is instead
$\Gamma_\star/\mathcal C_\beta(L)
=\Gamma_\star/[2L\Lambda(\pi L/\beta)]$.  At the de Sitter temperature this
is $\Gamma_\star/[R C_{\rm opt}(L/R)]$.

\subsection{Smooth compact sources}

The hard endpoint in \cref{eq:linear-profile} is useful for exact constants
but is not a smooth source.  We now show that it is a legitimate limiting
profile and that every strict window persists.

\begin{lemma}[Smooth density closure]
\label{lem:smooth-density}
Fix an optical source radius $\ell>0$ and a normalized
$p\in L^2(0,\ell)$.  There are normalized
$p_n\in C_c^\infty([0,\ell))$, regular at the origin, such that
\begin{equation}
 p_n\longrightarrow p
 \quad\hbox{in }L^2(\R_+).
 \label{eq:smooth-convergence}
\end{equation}
Their thermal dephasing forms converge, and their cumulative energy measures
converge uniformly.
\end{lemma}

\begin{proof}
Density of smooth functions supported strictly inside $(0,\ell)$ gives
\cref{eq:smooth-convergence}; such functions are automatically regular radial
data at the origin.  On data supported in $[0,\ell]$,
\begin{equation}
 P_\ell h_0^{-1}\coth(\beta h_0/2)P_\ell
 \leq P_\ell h_0^{-1}P_\ell
 +\frac2\beta P_\ell h_0^{-2}P_\ell
 \label{eq:thermal-compressed-bound}
\end{equation}
as quadratic forms.  The right side is bounded by
\cref{lem:compressed-inverses}; hence the covariance is continuous under
$L^2$ convergence.  Finally,
\begin{equation}
 \sup_{0\leq x\leq\ell}
 \left|\int_0^x(p_n^2-p^2)\dd s\right|
 \leq\norm{p_n-p}_2(\norm{p_n}_2+\norm p_2),
 \label{eq:mass-uniform-convergence}
\end{equation}
which proves uniform convergence of cumulative energy.
\end{proof}

It remains to produce a supported final datum by a smooth spacetime source.
This also proves sharpness over sources when the source ball approaches the
declared final-support ball; it does not prove exact controllability from a
strictly smaller cylinder of optical radius $L-T$.  Let
$\phi_{\rm free,n}$ be the homogeneous solution with final data
$(q,p)=(0,p_n)$.  Choose a smooth time cutoff $\eta_n$ that vanishes before
the protocol, equals one near the final slice, and changes over a time shorter
than the optical gap between $\supp p_n$ and $\ell$.  If $\mathcal P$ is the
conformal wave operator, set
\begin{equation}
 J_n=\mathcal P(\eta_n\phi_{\rm free,n}).
 \label{eq:smooth-source}
\end{equation}
The retarded solution sourced by $J_n$ is
$\eta_n\phi_{\rm free,n}$.  It therefore has exactly the desired final data.
Finite propagation keeps $J_n$ inside the declared optical worldtube.  This
construction proves smooth source existence, but it does not bound the stress
of an autonomous mechanism that generates the prescribed $J_n$.

For a finite pointer, apply this construction to each conditional datum
separately and take a common cutoff interval and worldtube.  Because only
finitely many branches are present, the approximations can be chosen so that
all branch energies and all pairwise thermal forms converge simultaneously.
The resulting diagonal controlled source therefore realizes the Schur
channel used in \cref{thm:finite-pointer-renyi}.  This observation establishes
finite-pointer source closure; it does not add a general-$d$ saturation claim.
